\section{正交基、Gram--Schmidt 与 QR 分解}
\label{sec:03-Linear-Algebra/03-Orthogonality/04}

假设你已经有两个基向量

\[
a_1=\begin{pmatrix}1\\1\\1\end{pmatrix},\qquad
a_2=\begin{pmatrix}1\\0\\2\end{pmatrix},
\]

它们张成 \(\R^3\) 中的一个平面。现在手上有一个该平面里的向量，比如

\[
p=\begin{pmatrix}3\\1\\4\end{pmatrix}=1\cdot a_1+2\cdot a_2,
\]

你想知道它在基 \(\{a_1,a_2\}\) 下的坐标——也就是那两个系数 \(1\) 和 \(2\)。

如果只有 \(p\) 而不知道系数，你得解方程组：

\[
c_1\begin{pmatrix}1\\1\\1\end{pmatrix}+c_2\begin{pmatrix}1\\0\\2\end{pmatrix}
=\begin{pmatrix}3\\1\\4\end{pmatrix}.
\]

写成矩阵：

\[
\begin{pmatrix}1&1\\1&0\\1&2\end{pmatrix}
\begin{pmatrix}c_1\\c_2\end{pmatrix}
=\begin{pmatrix}3\\1\\4\end{pmatrix}.
\]

这是个超定系统，但既然 \(p\) 确实在列空间里，它应该有精确解——实际上解出来就是 \(c_1=1,c_2=2\)。然而这里的困难不在于有没有解，而在于求解过程：两个基向量不是正交的（\(a_1^{\trans}a_2=3\neq0\)），所以从 \(p\) 中提取 \(c_1\) 时，\(a_2\) 的分量也会混进投影里去。你不能简单地说“\(c_1\) 等于 \(p\) 和 \(a_1\) 的内积除以长度平方”——那样算出来的是 \(p\) 在 \(a_1\) 方向上的投影长度，不是 \(a_1\)-\(a_2\) 混合表示中的坐标。

普通基能表示子空间，但坐标提取需要解方程。正交基消除了这种耦合。

\subsection{正交归一基让坐标随手可得}

设 \(Q\) 的列 \(q_1,\ldots,q_n\) 两两正交且各自归一：

\[
q_i^{\trans}q_j=
\begin{cases}
1,&i=j,\\
0,&i\neq j.
\end{cases}
\]

写成矩阵就是

\[
Q^{\trans}Q=I.
\]

现在对任何列空间中的向量 \(p=Qx\)，想知道坐标 \(x\)。两边左乘 \(Q^{\trans}\)：

\[
Q^{\trans}p = Q^{\trans}Qx = Ix = x.
\]

这就是正交基的威力：要拿到第 \(i\) 个坐标，只需计算 \(q_i^{\trans}p\)。没有联立方程，不需要求逆，一个内积就够了。

投影矩阵也随之简化。一般投影要用

\[
A(A^{\trans}A)^{-1}A^{\trans},
\]

但如果列已经是正交归一的，这个公式直接塌缩为

\[
P=QQ^{\trans}.
\]

所以一个很实际的目标自然浮现：如果手中的基不正交，能不能把它变成一个张成同一子空间的正交归一基？Gram--Schmidt 正交化就是要做这件事。

\subsection{“正交矩阵”与“正交归一列”的区分}

先说清楚一个术语上的坑。矩形矩阵也可以满足 \(Q^{\trans}Q=I\)——这表示列的几何性质很好，但列数小于行数，\(Q\) 不是方阵。这时候我们称它的列正交归一，而不是称它为“正交矩阵”。

“正交矩阵”这个叫法保留给方阵的情形。当 \(Q\) 是 \(n\times n\) 方阵且 \(Q^{\trans}Q=I\) 时，左逆等于右逆：

\[
Q^{-1}=Q^{\trans},
\]

并且反过来也有 \(QQ^{\trans}=I\)。这类矩阵同时保持长度和夹角：

\[
\|Qx\|_2^2 = x^{\trans}Q^{\trans}Qx = \|x\|_2^2,
\qquad
(Qx)^{\trans}(Qy) = x^{\trans}Q^{\trans}Qy = x^{\trans}y.
\]

旋转和镜像是两个典型。矩形 \(Q\) 不会把整个 \(\R^m\) 变换回自身——行数多出来，像空间只占一个低维子空间。但它仍然把低维坐标 \(x\) 等距地嵌入列空间：\(\|Qx\|_2=\|x\|_2\)。换句话说，从 \(x\) 到 \(Qx\)，长度不变，只是在更高维的空间里换了一个姿势。

\subsection{Gram--Schmidt：不断减去已经解释的方向}

现在动手把任意一组线性无关的向量 \(a_1,\ldots,a_n\) 变成正交归一组 \(q_1,\ldots,q_n\)。

先从 \(a_1\) 开始。还没有任何旧方向需要避开，直接归一化：

\[
u_1=a_1,\qquad q_1=\frac{u_1}{\|u_1\|_2}.
\]

轮到 \(a_2\)。它和 \(q_1\) 方向不是正交的——\(a_2\) 中有一部分分量躺在 \(q_1\) 的方向上。这个分量可以在 \(q_1\) 上投影出来：

\[
(q_1^{\trans}a_2)q_1.
\]

把 \(a_2\) 中已经可以被 \(q_1\) 解释的部分削掉，剩下的就是新的、与 \(q_1\) 正交的方向：

\[
u_2 = a_2 - (q_1^{\trans}a_2)q_1.
\]

再归一化：

\[
q_2 = \frac{u_2}{\|u_2\|_2}.
\]

这样得到了两个正交归一的向量，它们和原来的 \(a_1,a_2\) 张成同一个平面。

一般步骤：假设已经构造出 \(q_1,\ldots,q_{k-1}\)，它们两两正交归一，且张成 \(\operatorname{span}\{a_1,\ldots,a_{k-1}\}\)。处理 \(a_k\) 时，先算出 \(a_k\) 在前 \(k-1\) 个方向上的所有投影，一并削去：

\[
u_k = a_k - \sum_{j=1}^{k-1}(q_j^{\trans}a_k)q_j,
\qquad
q_k = \frac{u_k}{\|u_k\|_2}.
\]

每一步的减法都是一句相同的逻辑：从新向量里，把已经被已有子空间解释掉的部分剔干净，只保留还在子空间之外的新方向。因为每次只减去了前面向量的线性组合，前 \(k\) 个 \(q\) 和前 \(k\) 个 \(a\) 始终张成同一个子空间。

如果某一步算出 \(u_k=0\)，说明 \(a_k\) 已经在前 \(k-1\) 个向量的张成里——输入向量组不是线性无关的。这时候归一化无法进行。这不是程序 bug，是算法忠实地告诉你：这一列没有贡献新方向。失败的位置本身就是信息。

\subsection{一个完整例子}

回到开头的两个向量：

\[
a_1=\begin{pmatrix}1\\1\\1\end{pmatrix},\qquad
a_2=\begin{pmatrix}1\\0\\2\end{pmatrix}.
\]

第一步：

\[
\|a_1\|_2=\sqrt{1^2+1^2+1^2}=\sqrt3,
\qquad
q_1=\frac1{\sqrt3}\begin{pmatrix}1\\1\\1\end{pmatrix}.
\]

第二步，先问：\(a_2\) 在 \(q_1\) 方向上有多少？

\[
q_1^{\trans}a_2 = \frac{1\cdot1+1\cdot0+1\cdot2}{\sqrt3} = \frac{3}{\sqrt3} = \sqrt3.
\]

所以

\[
(q_1^{\trans}a_2)q_1 = \sqrt3\cdot\frac1{\sqrt3}\begin{pmatrix}1\\1\\1\end{pmatrix}
= \begin{pmatrix}1\\1\\1\end{pmatrix}.
\]

从 \(a_2\) 中减去它：

\[
u_2 = \begin{pmatrix}1\\0\\2\end{pmatrix} - \begin{pmatrix}1\\1\\1\end{pmatrix}
= \begin{pmatrix}0\\-1\\1\end{pmatrix}.
\]

归一化：

\[
\|u_2\|_2=\sqrt{0^2+(-1)^2+1^2}=\sqrt2,\qquad
q_2=\frac1{\sqrt2}\begin{pmatrix}0\\-1\\1\end{pmatrix}.
\]

现在验证：\(q_1^{\trans}q_2 = \frac{1}{\sqrt3\cdot\sqrt2}(1\cdot0+1\cdot(-1)+1\cdot1)=0\)，两者的长度都是 \(1\)。用 \(q_1,q_2\) 张成的平面和原来的 \(a_1,a_2\) 完全一样。

关键是第二步中那个问题：“\(a_2\) 中有多少已经被 \(a_1\) 的方向解释？”减掉答案后剩下的，才是新信息。

\subsection{把过程写成矩阵分解}

每一步都把 \(a_k\) 写成 \(q_1,\ldots,q_k\) 的组合：

\[
a_k = r_{1k}q_1 + r_{2k}q_2 + \cdots + r_{kk}q_k.
\]

不会出现 \(q_{k+1},q_{k+2},\ldots\)，因为后面那些方向还没被生成，\(a_k\) 也还没来得及在它们上面留下分量。所以系数矩阵 \(R\) 是上三角的。把所有列合在一起：

\[
A = QR,
\quad\text{其中 }A=\begin{pmatrix}a_1&\cdots&a_n\end{pmatrix},\;
Q=\begin{pmatrix}q_1&\cdots&q_n\end{pmatrix},\;
R=(r_{ij})_{n\times n}\text{ 是上三角矩阵}.
\]

因为 \(Q^{\trans}Q=I\)，直接用 \(Q^{\trans}\) 就能读出 \(R\)：

\[
Q^{\trans}A = Q^{\trans}QR = R.
\]

具体到前面的例子：

\[
Q=
\begin{pmatrix}
1/\sqrt3&0\\
1/\sqrt3&-1/\sqrt2\\
1/\sqrt3&1/\sqrt2
\end{pmatrix},
\qquad
R=
\begin{pmatrix}
\sqrt3&\sqrt3\\
0&\sqrt2
\end{pmatrix}.
\]

\(r_{11}=\|a_1\|=\sqrt3\)，\(r_{12}=q_1^{\trans}a_2=\sqrt3\)，\(r_{22}=\|u_2\|=\sqrt2\)，\(r_{21}=0\)。直接乘回去：

\[
QR=
\begin{pmatrix}
1/\sqrt3&0\\
1/\sqrt3&-1/\sqrt2\\
1/\sqrt3&1/\sqrt2
\end{pmatrix}
\begin{pmatrix}
\sqrt3&\sqrt3\\
0&\sqrt2
\end{pmatrix}
=
\begin{pmatrix}
1&1\\
1&0\\
1&2
\end{pmatrix}
=A.
\]

\(A=QR\) 可以看作与消元中的 \(A=LU\) 并列的另一种分解：\(LU\) 追求三角结构便于前代回代；QR 同时提供正交几何和三角计算结构。

\subsection{QR 怎样求最小二乘}

现在把 QR 装进最小二乘问题里。

设 \(A\) 满列秩，已有简约 QR 分解 \(A=QR\)，其中 \(Q\) 列正交归一，\(R\) 可逆上三角。最小二乘的目标是

\[
\min_x\|QRx-b\|_2.
\]

把 \(b\) 拆成两部分：落在列空间里的和落在正交补里的：

\[
b = QQ^{\trans}b + (I-QQ^{\trans})b.
\]

第一项是 \(b\) 在 \(\operatorname{Col}(A)=\operatorname{Col}(Q)\) 上的投影；第二项是与之正交的残差。代入：

\[
\begin{aligned}
QRx-b
&= QRx - QQ^{\trans}b - (I-QQ^{\trans})b \\
&= Q(Rx - Q^{\trans}b) - (I-QQ^{\trans})b.
\end{aligned}
\]

右边两项正交（第一项在列空间，第二项在正交补里），因此

\[
\|QRx-b\|_2^2
= \|Q(Rx - Q^{\trans}b)\|_2^2 + \|(I-QQ^{\trans})b\|_2^2.
\]

\(Q\) 保持长度，所以第一项就是 \(\|Rx - Q^{\trans}b\|_2^2\)。第二项和 \(x\) 无关——它是无论如何都无法消除的那部分残差，由 \(b\) 和列空间的距离决定。

要让第一项最小到零，只需解上三角方程

\[
R\hat x = Q^{\trans}b.
\]

没有显式形成 \(A^{\trans}A\)，没有求逆。先用 \(Q^{\trans}\) 提取 \(b\) 在列空间各正交方向上的坐标，再通过三角回代反解出参数。上一节末尾的警告——“形成 \(A^{\trans}A\) 会平方条件数”——正是从这里得到回应的：QR 避开了那个平方步骤，直接用更稳定的矩阵 \(R\) 代替了 \(A^{\trans}A\)。

\subsection{经典 Gram--Schmidt 与数值计算中的 QR}

经典 Gram--Schmidt 把正交化的思路写得明明白白：逐步减投影，逐步归一。但在浮点运算中，前面得到的方向会因为舍入误差而微微偏离完全正交的状态，后面的减法又基于这些已经偏了的方向，误差会累积。

实际数值软件用改进 Gram--Schmidt（调整投影顺序以减少误差传递）、Householder 反射或 Givens 旋转。它们追求的仍然是同一个结果 \(A=QR\)，但用更稳定的运算次序来达成。

因此要分清楚两件事：

\begin{itemize}
\tightlist
\item
  想理解 QR 从哪里来、正交化为什么是“减投影”——经典 Gram--Schmidt 是最透明的推导；
\item
  想在计算机上可靠地算出 QR——Householder 等方法才是正确的工具。
\end{itemize}

一个公式在精确算术里完全正确，不代表它就是浮点环境中的最优实现。这个区分在数值分析部分会反复出现。

\subsection{回看整条主线}

这一单元从一条最简单的几何事实出发：垂直时长度平方可以直接相加。由此引出正交、正交补和空间分解；最近点的残差必须落在正交补中，于是得到了投影；无解方程借助投影变成了最小二乘；为了让投影和最小二乘更容易计算，又需要正交基；Gram--Schmidt 把任意基变成正交基；QR 把正交基和三角结构封装成一个分解。

整条链条由同一个问题逐层推进，每个概念都是前一个的自然应答。

延伸阅读：MIT 18.06 Lecture 17：Orthogonal Matrices and Gram--Schmidt。
